\section{Discussion}
\label{sec:discussion}

\subsection{Main empirical takeaways}

Three findings are consistent in the current run family. First, speculative decoding
provides consistent acceleration across all tested configurations
($S>1$ throughout). Second, speedup declines as $k$ increases in this run
family, with the strongest performance at $k{=}4$ and the weakest at
$k{=}16$. Third, deterministic decoding outperforms stochastic decoding in 5 of
6 matched draft--$k$ pairs.

Together, these results suggest that practical deployment should prioritise a
smaller block size ($k{=}4$) and deterministic policy when latency is the
primary objective.

The key engineering implication is that, on consumer-class hardware, speedup
is governed by a balance between acceptance and per-cycle systems overhead.
In our runs, $B_{\text{eff}}$ grows with $k$ while speedup still declines,
indicating that memory-traffic and loop-overhead growth outpace useful
acceptance gains at larger block sizes.

The ACSD extension adds a second practical option: adaptive control keeps
speed close to best fixed-policy settings while providing explicit rescue and
tail-fallback mechanisms. On the current run family, ACSD reaches
$S{=}2.9792$ (deterministic) and $S{=}2.8981$ (stochastic), with low fallback
activation (1.6\% and 2.2\%).

From a positioning standpoint, this matters because demand for local and
consumer-hardware LLM inference is increasing: the relevant question is no
longer only ``how fast can we decode in data centers,'' but ``which decoding
policy remains efficient and stable on constrained personal hardware.''

For this reason, our recommendations are framed as hardware-profiled operating
points rather than universal defaults.

\subsection{Draft-size implications}

The 0.5B draft produces the best single configuration
(0.5B, $k{=}4$, deterministic), and also has a slightly higher mean speedup
than 1.5B across the full grid in this run family. This indicates that
draft-size selection is hardware- and budget-sensitive: the smaller draft can
win on end-to-end speed even when the larger draft improves acceptance.

\subsection{Speed versus quality}

Quality drift is task-dependent. CNN/DailyMail and MMLU remain relatively
stable across settings, while GSM8K is more variable, especially under
stochastic decoding. A conservative filter
($\max |\Delta| \le 1$ across all three tasks) retains a small subset of
configurations, among which 1.5B, $k{=}8$, deterministic remains best.
This supports a practical policy: use deterministic mode with $k{=}4$ for
production-like latency targets, and treat stochastic settings as
quality-sensitive configurations requiring additional calibration.
Final winner claims should be interpreted together with paired confidence
intervals and corrected pairwise tests from a coherent artifact family.

\section{Threats to validity}
\label{sec:threats}

\textbf{Artifact consistency across reruns.}
Some baseline files were regenerated after initial summary artifacts were
produced, which can create mismatches between raw baseline totals and
summary-implied baselines. To mitigate this, we use one coherent summary source
for comparative claims and report this limitation explicitly.

\textbf{Single-hardware scope.}
All results are from one RTX~4090 Laptop. Relative rankings are informative for this
deployment class, but absolute speedups may shift on other GPU generations.
In this run family, CUDA-graph metadata indicates pending/zero capture rate,
so reported gains reflect the same code path under fallback execution rather
than successful graph replay.

\textbf{Model-family scope.}
All pairings are same-family (Qwen2.5 target and drafts). Cross-family
pairings may produce materially different acceptance and quality behaviour.

\textbf{Evaluation breadth.}
The benchmark set spans math reasoning, knowledge QA, and summarisation, but
does not exhaust safety, factuality, or long-context behaviour. ROUGE-L in
particular is an incomplete summarisation quality proxy.

\section{Conclusion}
\label{sec:conclusion}

This paper provides a controlled consumer-GPU evaluation of speculative
decoding with a Qwen2.5-3B target and same-family drafts. The current evidence
shows consistent acceleration across all tested configurations, with best
observed performance at 0.5B, $k{=}4$, deterministic
($S{=}3.0674$). We find that increasing $k$ from 4 to 8 and 16 reduces net
speedup in this setup, and that deterministic decoding is generally faster than
stochastic decoding for matched settings.

These findings yield an actionable operating recommendation for the present
hardware budget: prefer deterministic speculative decoding at $k{=}4$, with
0.5B draft as the default fast path and larger drafts reserved for
quality-sensitive scenarios. The ACSD results (Section~\ref{sec:acsd}) show
that adaptive draft length and cascaded draft-model switching can deliver
near-best fixed-policy speed while adding low-frequency rescue/fallback
controls for tail robustness. This supports ACSD as a practical consumer-GPU
policy layer rather than a replacement for the fixed-policy fast path.
